% ==============================================================
% 第3章 系统需求分析
% ==============================================================
\chapter{系统需求分析}

\section{系统功能需求分析}

基于社区2型糖尿病慢性病管理的业务场景，本系统面向\textbf{患者}与\textbf{社区医生}（兼管理职能）两类用户群体。结合对社区卫生服务中心慢病管理工作的调研分析，系统需围绕健康数据采集、风险智能评估、随访任务调度、用药依从管理以及医生协同工作五个核心业务领域提供功能支撑，各模块的具体需求分析如下。

\subsection{健康数据采集需求}

社区2型糖尿病患者以中老年群体为主，其数字化工具使用能力参差不齐，传统纯文字录入方式存在一定的操作门槛，容易导致健康数据上报不及时、不完整等问题。为此，系统需以语音录入为主、文字录入为辅，允许患者以日常口语化表述完成健康数据的上报，同时保留文字录入途径作为补充，以兼顾不同用户的使用习惯。在录入数据类型方面，系统需支持血糖值（空腹/餐后分类）、血压（收缩压与舒张压）、体重等常规体征指标，并支持患者记录每日用药打卡情况与情绪状态，作为健康管理的辅助参考信息。考虑到语音识别存在误解析的可能，系统还需在数据提交前向患者展示解析后的结构化结果，由患者确认无误后方可入库，以保障录入数据的准确性与可靠性。

\subsection{风险评估与预警需求}

社区医生需统筹管理辖区内数量较多的在管患者，若依赖人工逐一审阅健康数据，不仅效率低下，且易因疏漏而延误对高风险患者的及时干预。为解决这一问题，系统需能够在患者提交健康数据后自动完成健康状态评估，依据2型糖尿病临床管理相关标准，综合患者的血糖控制水平、血糖波动情况及异常持续时长等维度，输出红、黄、绿三级风险预警结果，无需医生手动触发。对于处于红码预警状态的高风险患者，系统需向责任医生实时推送通知，确保异常情况得到及时响应。此外，系统还需为每位患者动态维护健康画像，直观呈现近期血糖达标率、波动趋势等关键指标，供医生快速全面地掌握患者健康状态。

\subsection{随访调度需求}

现有社区随访工作多依赖人工排期，随访周期固定，难以响应患者病情的动态变化，存在高风险患者未能及时跟进的管理盲区。为此，系统需能够根据患者当前的风险等级自动调整其随访周期，对病情不稳定的患者缩短随访间隔，对状态平稳的患者则适当延长，实现随访资源的合理分配。在随访形式上，系统需支持线上轻问诊、线下门诊及上门巡诊等多种方式，并结合患者风险程度与当前就诊资源进行合理分流推荐。随访任务生成后，系统应自动通知责任医生，并在医生工作台以结构化列表形式展示待办事项，便于统一跟进管理。与此同时，患者端也需能查看本人的下次随访安排，以提升患者对慢病管理的参与感与配合度。

\subsection{用药管理需求}

用药依从性差是2型糖尿病患者血糖控制不达标的重要原因之一，患者因遗忘服药或药品断供而导致病情波动的情况在社区管理中较为普遍。针对这一问题，系统需根据患者的处方信息在约定时间推送服药提醒，并支持同一患者配置多个服药时间点，以适应胰岛素注射与口服降糖药并用的复杂用药场景。在依从性管理方面，系统需根据患者的日常打卡记录自动统计其用药依从率，并以可视化形式呈现于医生工作台，辅助医生评估患者的用药配合情况。此外，系统还需根据患者录入的剩余药量与每日用药剂量估算药品可用天数，在库存不足时提前向患者及责任医生发出续方提醒，从源头降低因断药引发病情加重的风险。

\subsection{医生协同工作台需求}

社区医生同时承担患者日常管理、随访跟进与基础行政事务等多重职责，亟需一个能够整合患者信息与工作任务的统一操作平台。为此，系统需为医生提供患者总览看板，以列表形式集中展示辖区内全部在管患者，并支持按风险等级进行筛选与排序，优先呈现高风险患者，便于医生快速识别重点关注对象。在预警管理方面，系统需实时汇聚展示当前所有红码患者的异常信息，支持医生快速调阅对应患者的历史健康记录，辅助临床判断。在任务管理方面，系统需将自动生成的随访任务以结构化形式呈现于工作台，医生可查看任务详情并更新任务处理状态（待处理/已联系/已完成），实现随访工作的全程可追踪。此外，系统还需支持医生进行患者档案管理，包括新增在管患者及维护患者基础信息（如姓名、年龄、病程年限、用药方案等），兼履行系统管理员职能，无需另行设置独立的管理员角色。


\subsection{知识辅助服务需求}

在临床实践中，患者对自身体征数据的含义理解不足，医生在面对大量在管患者时需快速获取基于指南的诊疗参考，这两类需求均无法通过简单的数据展示或固定模板满足。为此，系统需具备基于权威临床知识的智能辅助能力：在患者端，系统应在患者提交健康数据并完成风险评估后，自动生成基于《中国糖尿病防治指南》等权威来源的健康反馈建议，以通俗易懂的语言向患者解读当前体征状况并给出生活管理指导，所有建议均需附加免责提示以明确系统的辅助工具定位；在医生端，系统应在医生查看患者详情时，综合该患者的风险等级、异常指标、近期体征趋势与用药依从率等多维数据，自动生成结构化的诊疗辅助摘要，引用相关指南条目为医生临床决策提供参考依据。上述知识辅助功能应基于检索增强生成（RAG）技术实现，确保输出内容有据可查、可追溯至具体知识来源。


% ------------------------------------------------------------------
\section{系统非功能需求分析}

\textcolor{red}{[TODO: 补充非功能需求分析，包括以下维度]}

\subsection{性能需求}

系统在常规使用条件下，页面响应时间应控制在3秒以内。服务端ASR语音转写延迟不超过3秒（基于15秒以内的录音片段）。LLM语音解析响应时间不超过5秒。TTS语音合成延迟不超过3秒（基于200字以内的文本）。风险评估计算应在数据提交后1秒内完成。系统应支持同时在线用户数不低于100人。

\subsection{安全性需求}

系统应对用户登录密码进行加密存储，采用哈希算法（如bcrypt）保障密码安全。系统应实现基于角色的访问控制（RBAC），确保患者仅能访问本人的健康数据，医生仅能访问其管辖患者的数据。系统应对敏感医疗数据的传输采用HTTPS加密协议。

\subsection{可用性需求}

系统界面应遵循适老化设计原则，采用大字体（不小于16px）、高对比度配色方案。核心操作流程应控制在3步以内完成。语音交互功能应提供清晰的操作引导和反馈提示。系统应支持将健康反馈建议以语音播报形式呈现，降低老年患者的阅读负担。

\subsection{可维护性与可扩展性需求}

系统采用模块化设计，各功能模块间应保持低耦合。Agent模块应支持独立更新和替换，不影响其他模块正常运行。LLM调用层应提供统一封装接口，支持底层模型的灵活切换。语音识别（ASR）与语音合成（TTS）服务应采用策略模式封装，支持通过配置文件切换底层供应商（如Whisper/火山引擎/通义千问），实现供应商无关性。知识库应支持在线更新，无需重启系统。系统所有AI交互过程（语音解析、健康反馈、诊疗摘要、风险评估）的输入与输出均应持久化存储至Agent交互日志表，支持审计回查与历史追溯，确保AI辅助决策的可追溯性与可解释性。


% ------------------------------------------------------------------
\section{系统用例分析}

\subsection{系统参与者识别}

根据功能需求分析，本系统识别出以下两类主要参与者（Actor）：

\begin{itemize}
  \item \textbf{患者}：社区在管的2型糖尿病患者，系统的核心服务对象。患者通过系统录入健康数据、查看风险评估结果、接收用药提醒、查看随访安排。
  \item \textbf{社区医生}：负责辖区内糖尿病患者管理的基层医生，同时兼任系统管理员角色。医生通过系统查看患者风险预警、管理随访任务、维护患者档案。
\end{itemize}

\subsection{系统总体用例图}

\todofig{}

\textcolor{red}{[TODO: 插入系统总体用例图，使用draw.io绘制后导出]}

图~\ref{fig:usecase-overall}展示了系统的总体用例图，描述了患者和社区医生与系统各功能模块之间的交互关系。

\subsection{核心用例描述}

\subsubsection{用例：语音录入健康数据}

\begin{table}[H]
\centering
\caption{用例描述：语音录入健康数据}
\label{tab:usecase-voice}
\begin{tabularx}{\textwidth}{|l|X|}
\hline
\textbf{用例名称} & 语音录入健康数据 \\
\hline
\textbf{参与者} & 患者 \\
\hline
\textbf{前置条件} & 患者已登录系统 \\
\hline
\textbf{后置条件} & 健康数据成功存入系统，触发风险评估流程 \\
\hline
\textbf{基本流程} &
  1. 患者进入健康数据录入页面，选择语音录入模式 \newline
  2. 患者按住录音按钮，口述当日健康数据 \newline
  3. 录音结束后，系统将音频上传至服务端，调用ASR服务转写为文本 \newline
  4. 患者点击"AI解析"，系统调用LLM将口语化文本解析为结构化数据 \newline
  5. 系统向患者展示解析结果，请求确认 \newline
  6. 患者确认数据无误 \newline
  7. 系统保存数据并自动触发风险评估，同时将解析过程记录至Agent交互日志 \\
\hline
\textbf{备选流程} &
  3a. ASR转写结果有误，患者可手动编辑文本后再点击"AI解析" \newline
  5a. 结构化解析结果有误，患者选择手动修正后提交 \newline
  5b. 患者放弃语音方式，切换为文字录入 \\
\hline
\end{tabularx}
\end{table}

\subsubsection{用例：查看患者风险预警}

\begin{table}[H]
\centering
\caption{用例描述：查看患者风险预警}
\label{tab:usecase-risk}
\begin{tabularx}{\textwidth}{|l|X|}
\hline
\textbf{用例名称} & 查看患者风险预警 \\
\hline
\textbf{参与者} & 社区医生 \\
\hline
\textbf{前置条件} & 医生已登录系统；存在新的风险预警记录 \\
\hline
\textbf{后置条件} & 医生已知晓高风险患者情况 \\
\hline
\textbf{基本流程} &
  1. 医生登录后进入工作台主页 \newline
  2. 系统在预警区域展示当前红码患者列表 \newline
  3. 医生点击某患者，查看异常指标详情和历史健康记录 \newline
  4. 医生根据情况决定后续处理（联系患者/安排随访/转诊） \\
\hline
\textbf{备选流程} &
  2a. 当前无红码预警，显示``暂无紧急预警'' \\
\hline
\end{tabularx}
\end{table}

\subsubsection{用例：管理随访任务}

\begin{table}[H]
\centering
\caption{用例描述：管理随访任务}
\label{tab:usecase-visit}
\begin{tabularx}{\textwidth}{|l|X|}
\hline
\textbf{用例名称} & 管理随访任务 \\
\hline
\textbf{参与者} & 社区医生 \\
\hline
\textbf{前置条件} & 系统已根据风险等级自动生成随访任务 \\
\hline
\textbf{后置条件} & 随访任务状态已更新 \\
\hline
\textbf{基本流程} &
  1. 医生在工作台查看待办随访任务列表 \newline
  2. 医生选择某一任务，查看患者信息和任务详情 \newline
  3. 医生完成随访后，更新任务状态为``已完成'' \newline
  4. 系统记录完成时间，从待办列表中移除 \\
\hline
\textbf{备选流程} &
  3a. 患者无法联系，医生将任务状态更新为``已延期''并备注原因 \\
\hline
\end{tabularx}
\end{table}

\subsubsection{用例：查看健康反馈建议}

\begin{table}[H]
\centering
\caption{用例描述：查看健康反馈建议}
\label{tab:usecase-feedback}
\begin{tabularx}{\textwidth}{|l|X|}
\hline
\textbf{用例名称} & 查看健康反馈建议 \\
\hline
\textbf{参与者} & 患者 \\
\hline
\textbf{前置条件} & 患者已完成一次健康数据录入并获得风险评估结果 \\
\hline
\textbf{后置条件} & 患者查看到基于指南的个性化健康反馈建议 \\
\hline
\textbf{基本流程} &
  1. 患者提交健康数据后，系统完成风险评估 \newline
  2. 系统自动根据本次体征数据与风险等级，从知识库检索相关指南建议 \newline
  3. 系统生成通俗易懂的健康反馈文本，展示于评估结果页面 \newline
  4. 患者阅读反馈建议，了解当前健康状况及改善方向 \\
\hline
\textbf{备选流程} &
  3a. 知识库无匹配内容或生成服务不可用，页面仅展示风险评估结果，不展示反馈建议 \\
\hline
\end{tabularx}
\end{table}

\textcolor{red}{[TODO: 根据需要补充更多用例描述，如：用药打卡、续方提醒、患者建档等]}
